\documentclass{axiobac}

\annee{1990}
\mois{Juin}
\session{principale}
\epreuve{Mathématiques}
\sectionbac{Math-Sciences \& Math-Technique}
\sections{mathematiques, sciences-techniques}
\duree{4 heures}
\coefficient{5}
\ministere{Ministère de l'Éducation, de l'Enseignement Supérieur et de la Recherche Scientifique}

\begin{document}
\entete

%----------------------------------------------------------------------------------------
%   EXERCICE 1
%----------------------------------------------------------------------------------------
\begin{exercice}{4}
\titre{Équation du troisième degré et losange}
\chapitres{Nombres Complexes}

Pour tout nombre complexe $z$ on pose :
\[
\begin{aligned}
f(z) = {} & z^3 + 2\left(-\sqrt{3} + i\right) z^2 \\
          & + 4\left(1 - i\sqrt{3}\right) z + 8i .
\end{aligned}
\]

\begin{enumerate}
    \item
    \begin{enumerate}
        \item Montrer que l'équation $f(z) = 0$ possède une racine imaginaire pure $z_0$ que l'on déterminera.
        \item Résoudre alors l'équation $f(z) = 0$. \\
        On notera $z_1$ et $z_2$ les deux autres racines, $z_1$ étant celle qui a une partie imaginaire négative.
    \end{enumerate}
    \item On pose $\omega = \dfrac{z_1}{z_0}$.
    \begin{enumerate}
        \item Donner la forme trigonométrique de $\omega$.
        \item Le plan complexe est rapporté à un repère orthonormé direct $(O, \vec{u}, \vec{v})$. \\
        À tout nombre complexe $z$ non nul on associe les points $M$, $M_1$ et $M_2$ d'affixes respectives $z$, $\omega z$ et $\omega^2 z$. \\
        Montrer que $OMM_1M_2$ est un losange.
    \end{enumerate}
\end{enumerate}

\end{exercice}

%----------------------------------------------------------------------------------------
%   EXERCICE 2
%----------------------------------------------------------------------------------------
\begin{exercice}{4}
\titre{Rotation, antidéplacement et lieu de milieux}
\chapitres{Déplacements - Antidéplacements}

Dans le plan orienté, on considère un triangle $ABC$ non isocèle tel que $\left(\widehat{\overrightarrow{AB}, \overrightarrow{AC}}\right) \equiv \dfrac{\pi}{2} \ [2\pi]$. \\
À tout point $M$ de la droite $(AB)$ on associe le point $N$ de la droite $(AC)$ tel que $M$ et $N$ soient dans un même demi-plan de bord $(BC)$ et $BM = CN$.

\begin{enumerate}
    \item Montrer qu'il existe une unique rotation $\mathcal{R}$ telle que, pour tout point $M$ de $(AB)$, on a : $\mathcal{R}(M) = N$ et $\mathcal{R}(B) = C$. \\
    Préciser une mesure de son angle et construire son centre $\Omega$.
    \item Soit $O$ le milieu du segment $[BC]$. On désigne par $S_{(O\Omega)}$ la symétrie orthogonale d'axe $(O\Omega)$ et on pose : $f = S_{(O\Omega)} \circ \mathcal{R}$.
    \begin{enumerate}
        \item Déterminer $f(B)$ et $f(\Omega)$.
        \item Préciser la nature et les éléments caractéristiques de $f$.
    \end{enumerate}
    \item Soit $I$ le milieu du segment $[MN]$.
    \begin{enumerate}
        \item Quel est l'ensemble $D$ des points $I$ lorsque $M$ décrit la droite $(AB)$ ?
        \item Construire $D$.
    \end{enumerate}
\end{enumerate}

\end{exercice}

%----------------------------------------------------------------------------------------
%   PROBLÈME
%----------------------------------------------------------------------------------------
\begin{probleme}{12}

\partie{I}
\titre{Encadrer $\ln(n!)$ par des intégrales}
\chapitres{Intégrales, Fonction logarithme népérien, Suites réelles}

Soit la fonction
\[
f : \ ]-1, 1[ \longrightarrow \mathbb{R}, \qquad x \longmapsto f(x) = \int_0^x \frac{t^2}{1 - t^2}\,dt .
\]

\begin{enumerate}
    \item
    \begin{enumerate}
        \item Justifier l'existence de $f$.
        \item Montrer qu'il existe trois réels $\alpha$, $\beta$, $\gamma$ tels que :
        \[
        \forall t \in \mathbb{R} \setminus \{-1, 1\}, \qquad \frac{t^2}{1 - t^2} = \alpha + \frac{\beta}{1 - t} + \frac{\gamma}{1 + t} .
        \]
        \item En déduire que : $\forall x \in ]-1, 1[$, $f(x) = \dfrac{1}{2} \operatorname{Log}\left(\dfrac{1 + x}{1 - x}\right) - x$.
    \end{enumerate}
    \item Étudier les variations de $f$ et construire sa courbe représentative dans un repère orthonormé $(O, \vec{i}, \vec{j})$.
    \item
    \begin{enumerate}
        \item Montrer que : $\forall x \in \mathbb{R}^*_+$, $\forall k \in \mathbb{N}^*$, $\operatorname{Log} x \leqslant \dfrac{x}{k} - 1 + \operatorname{Log} k$.
        \item En déduire que :
        \[
        \forall k \in \mathbb{N}^*, \quad \int_{k - \frac12}^{k + \frac12} \operatorname{Log} x\,dx \leqslant \operatorname{Log} k
        \]
        et par suite :
        \[
        \forall n \in \mathbb{N}^*, \quad \int_{\frac12}^{n + \frac12} \operatorname{Log} x\,dx \leqslant \operatorname{Log}(n!) .
        \]
        \item Montrer que : $\forall n \in \mathbb{N}^*$, $\operatorname{Log}(n!) \geqslant \left(n + \frac12\right) \operatorname{Log}\left(n + \frac12\right) - n + \frac12 \operatorname{Log} 2$.
    \end{enumerate}
    \item Soit $(u_n)$ la suite définie sur $\mathbb{N}^*$ par : $u_n = \operatorname{Log}(n!) - \left(n + \frac12\right) \operatorname{Log} n + n$.
    \begin{enumerate}
        \item Montrer que : $\forall n \in \mathbb{N}^*$, $u_n \geqslant \frac12 \operatorname{Log} 2$.
        \item Vérifier que : $\forall n \in \mathbb{N}^*$, $u_n - u_{n+1} = (2n + 1)\, f\!\left(\dfrac{1}{2n + 1}\right)$.
        \item En déduire que $(u_n)$ est convergente.
    \end{enumerate}
\end{enumerate}

\partie{II}
\titre{Les intégrales $\int_0^1 x^{n/2}\sqrt{1-x}\,dx$}
\chapitres{Intégrales, Suites réelles}

Soit la suite $(v_n)$ définie par :
\[
v_0 = \int_0^1 (1 - x)^{\frac12}\,dx \qquad \text{et} \qquad \forall n \in \mathbb{N}^*, \qquad v_n = \int_0^1 x^{\frac{n}{2}} (1 - x)^{\frac12}\,dx .
\]

\begin{enumerate}
    \item Calculer $v_0$.
    \item
    \begin{enumerate}
        \item Le plan étant muni du repère orthonormé $(O, \vec{i}, \vec{j})$, déterminer l'ensemble des points $M(x, y)$ tels que : $y^2 - x(1 - x) = 0$.
        \item En déduire que $v_1 = \dfrac{\pi}{8}$.
    \end{enumerate}
    \item
    \begin{enumerate}
        \item Montrer que la suite $(v_n)$ est décroissante.
        \item En déduire qu'elle est convergente.
    \end{enumerate}
    \item
    \begin{enumerate}
        \item Prouver, à l'aide d'une intégration par parties, que : $\forall n \in \mathbb{N}$, $v_{n+2} = \dfrac{n + 2}{n + 5}\,v_n$.
        \item En déduire que :
        \[
        \forall n \in \mathbb{N}, \quad \frac{n + 2}{n + 5} \leqslant \frac{v_{n+1}}{v_n} \leqslant 1 \qquad \text{et} \qquad \lim_{n \to +\infty} \frac{v_{n+1}}{v_n} = 1 .
        \]
    \end{enumerate}
    \item
    \begin{enumerate}
        \item Montrer, par récurrence, que : $\forall n \in \mathbb{N}$, $v_n \cdot v_{n+1} = \dfrac{2\pi}{(n + 2)(n + 3)(n + 4)}$.
        \item En faisant apparaître, dans l'expression précédente, le rapport $\dfrac{v_{n+1}}{v_n}$, montrer que :
        \[
        \lim_{n \to +\infty} n^{\frac32}\, v_n = \sqrt{2\pi} .
        \]
    \end{enumerate}
    \item Montrer que : $\forall p \in \mathbb{N}$, $v_{2p} = \dfrac{2}{(2p + 1)(2p + 3)} \cdot \dfrac{\left(2^p \cdot p!\right)^2}{(2p)!}$.
\end{enumerate}

\partie{III}
\titre{La formule de Stirling : $\lim u_n = \ln\sqrt{2\pi}$}
\chapitres{Suites réelles, Fonctions exponentielles}

$(u_n)$ étant la suite définie dans I 4).

\begin{enumerate}
    \item Montrer que : $\forall n \in \mathbb{N}^*$, $e^{u_n} = \dfrac{n!}{n^{n + \frac12}} \cdot e^n$.
    \item Exprimer $e^{2u_p - u_{2p}}$ en fonction de $p$ et $v_{2p}$ ($p \in \mathbb{N}^*$).
    \item Soit $\ell$ la limite de la suite $(u_n)$. Déduire de ce qui précède que : $\ell = \operatorname{Log}\sqrt{2\pi}$. \\
    (On admettra que $\displaystyle\lim_{p \to +\infty} u_{2p} = \lim_{p \to +\infty} u_p = \ell$.)
\end{enumerate}

\end{probleme}

\end{document}
